\section{Conclusions}
\label{sec:conclusion}

This paper studied the minimum weighted feedback arc set problem through its ordering interpretation on nonnegative weighted directed graphs. The main contribution is a reproducible computational methodology rather than a new approximation framework. It combines an engineered local-ratio construction with topological add-back, a complementary weighted seed, and an incumbent-protected SCC-local refinement stage. Within that design, LR-TA serves as a strong constructive seed, while IPSNS is the main refinement framework that improves the incumbent only when the global backward-weight objective decreases.

The computational evidence supports a clear scoped conclusion. On the standard nonnegative sparse benchmark, IPSNS attains the best observed backward weight among the tested methods on 96 of 97 instances and improves on all tested internal and external comparators, including the strongest external baseline. On the full internal benchmark, the refinement stage never worsens the better internal seed, which confirms the intended non-degradation behavior in practice.

The exact-validation subset shows that the framework is near-optimal when bitmask dynamic programming certifies optima. The ablation study shows that the topological add-back phase is the dominant design element, with SCC-local refinement providing an additional but smaller gain.

The results also define an important boundary. On dense complete LOLIB ordering instances, the matrix-based pairwise-ordering baseline DRMacIver/FAS is stronger overall. This clarifies the regime in which the proposed framework is strongest: sparse nonnegative weighted digraphs. The method should be interpreted as a general weighted-digraph heuristic rather than as a dense-native linear-ordering solver.

The practical value of the framework lies in this combination of quality, structural focus, and reproducibility. The methodology is suitable for weighted ordering and cyclic-dependency repair settings where exact optimization is too expensive, sparse graph structure is meaningful, and refinement should not be allowed to damage a good seed. Within the stated boundaries, engineered local-ratio seeding plus incumbent-protected SCC refinement is an effective strategy for sparse weighted directed ordering problems.

Several directions follow from this evidence. The first is broader evaluation on additional sparse weighted digraph families, especially application-derived instances with heterogeneous edge-weight semantics. The second is stronger dense-native adaptation: the same incumbent-protected refinement discipline could be combined with a matrix-based or tournament-aware seed to extend competitive performance to dense complete ordering instances. The third is parallel or asynchronous SCC refinement, exploiting the structural independence of disjoint SCC neighborhoods to improve scalability on large sparse instances. The fourth is exact or ILP-based polishing on medium-sized instances, to close the quality reference gap between the small exact-validation subset and the full sparse benchmark.
